\chapter{Conclusion}
\label{ch:conclusion}

\section{Summary of Findings}

This thesis presented a reinforcement learning framework for articulated autonomous vehicle control, building on the foundational work from ME~780 (DDPG and Decision Transformer for single-vehicle path following) and ECE~682 (LQR control of the vehicle bicycle model with tractor-trailer kinematics). The key contributions demonstrated are:
\begin{itemize}
    \item a combined dynamic tractor and kinematic trailer model implemented in a custom Gymnasium simulation environment, extending the single-vehicle bicycle model from prior work;
    \item a vision-based observation space incorporating BEV images rendered from the Pygame simulation, enabling spatial perception of lane boundaries and obstacles;
    \item a TD3 reinforcement learning agent trained on lane-following and obstacle avoidance tasks with a multiplicative reward function penalising cross-track error, heading error, and hitch articulation angle;
    \item comparison of TD3 performance against PID and MPC classical control baselines and the prior ME~780 DDPG results.
\end{itemize}

The integration of visual perception via BEV imagery extends the state-only approach of the ME~780 work. By processing rendered BEV images through a CNN feature extractor, the agent gains spatial context about lane geometry and obstacle positions that vector state observations alone cannot provide.

\section{Limitations}

The study is subject to the following constraints:
\begin{itemize}
    \item Experiments are conducted entirely within the custom 2D Pygame/Gymnasium simulation environment; performance in higher-fidelity simulators or on physical hardware has not yet been characterised.
    \item The 2D simulation approximates the vehicle dynamics but does not capture all effects present in a real tractor-trailer system, such as load transfer, tire saturation at high slip angles, or hitch joint compliance.
    \item The current framework assumes a fixed payload; hitch dynamics change significantly with varying trailer mass.
    \item Obstacle avoidance is limited to static obstacles; dynamic obstacle handling is not addressed.
\end{itemize}

\section{Future Work}

Several directions remain open for future investigation:
\begin{itemize}
    \item \textbf{Higher-fidelity simulation}: Validating the trained policies in a physics-based 3D simulator (e.g., Gazebo) with a detailed URDF model of the tractor-trailer system, including realistic tire and hitch dynamics.
    \item \textbf{Hardware deployment}: Transferring trained policies to a physical robotic platform, requiring sim-to-real transfer techniques such as domain randomisation and system identification.
    \item \textbf{Decision Transformer for tractor-trailer}: Applying the DT approach from ME~780 to the articulated vehicle setting, with training data generated from the TD3 expert policy and augmented with high-error recovery trajectories.
    \item \textbf{Multi-trailer systems}: Extending the kinematic model and RL policy to double or B-train configurations.
    \item \textbf{Online adaptation}: Updating the policy in real time to compensate for varying payloads and road conditions.
    \item \textbf{Control Barrier Function overlays}: Implementing CBFs as a hard safety layer that overrides the RL agent when a collision or jackknife is imminent.
    \item \textbf{Dynamic obstacles}: Extending the obstacle avoidance environment to include moving obstacles and integrating predictive collision avoidance.
\end{itemize}
